%
%\documentclass{article}
%\documentclass{seminar}
%
%\usepackage[T1]{fontenc}
%\usepackage[utf8]{inputenc}
%\usepackage[english]{babel}
%\usepackage{graphicx}
%\usepackage{graphpap}
%\usepackage[pdftex,
%    pdftitle={ECE 474: Micro-C OS/II Overview},
%    pdfauthor={Blake Hannaford, University of Washington},
%    pdflang={en-US},
%    pdfsubject={ECE 474 Embedded Microcomputer Systems},
%    pdfkeywords={embedded systems, microcontrollers, ECE 474},
%    hidelinks,
%    unicode]{hyperref}
%
%\newpagestyle{BHstyle}{\sl U. of Washington \hfil EE 472 \hfil \thepage}{Hannaford \hfil \today}
%\pagestyle{headings}
%\slidepagestyle{BHstyle}
%\addtolength\oddsidemargin{3cm}
%
%\title{Micro-C OS/II Overview}
%
%
%
%\usepackage{fancyhdr}
%
%
%
%%%%%%%%%%%%%%%%%%%%%%%%%%%%%%%%%%%%%%%%%5
%%
%%  Set Up Margins
%\input{/users/home1/blake/templates/pagedim.tex}
%
%
%%%%%%%%%%%%%%%%%%%%%%%%%%%%%%%%%%%%%%%%%%%%%%%%%%
%%
%%         Page format Mods HERE
%%
%%Mod's to page size for this document
%\addtolength\textwidth{0cm}
%\addtolength\oddsidemargin{0cm}
%\addtolength\headsep{0cm}
%\addtolength\textheight{0cm}
%%\linespread{0.894}   % 0.894 = 6 lines per inch, 1 = "single",  1.6 = "double"
%
%
%
%%%%%%%%%%%%%%%%%%%%%%%%%%%%%% HEADER / FOOTER
%\pagestyle{fancy}
%%%%%%  Page header/footer fields for "NSF Style" proposal
%%%%%%  \section* will be changed with each section
%\lhead{\small\sc EE472}
%\rhead{}
\section*{\large\bf Introducing Micro-C OS/II ($\mu$Cos/II)}	%<chn>
%\lfoot{Hannaford, U. of Washington}
%\rfoot{Nov. 17, 2003}
%\cfoot{\thepage}
%%\renewcommand\headrulewidth{1pt}
%%\renewcommand\footrulewidth{1pt}
%
%
%\linespread{1.0}
%
	%<*>
%\begin{document}



%\begin{slide}
\section*{Introduction to MicroC/OS-II  }
%%%%** Section 1 
\section{Overview}
%%%%** Section 1.1
\subsection{Information Sources }
\begin{itemize}
\item http://www.uCOS-II.com/ \\
reference card\\
cool animations \\
Application notes \\
Examples of commercial products.
\item Book: ``MicroC/OS-II, The Real-Time Kernel,'' by
Jean J. Labrosse, CMP Books, 2nd edition, 2002.
\end{itemize}



%%%%** Section 1.2
\subsection{Characteristics and Features}

Ucos-II is a real time {\it kernel}.

``Operating System'' is more than a kernel.

Kernel supports
\begin{itemize}
  \item Scheduler
  \item Message passing (mailboxes)
  \item Synchronization (semphores)
  \item Memory Management
\end{itemize}

Does not support
\begin{itemize}
  \item I/O Devices
  \item File System
  \item Networking
\end{itemize}

%\end{slide}
%\begin{slide} 
Features

\begin{itemize}
  \item Source Code available (not ``GPL'').
  \item Free for educational use, license required
  for commercial use.
  \item Portable to multiple micro processors.
  \item ROMable.
  \item Scalable (only builds what you need)
  \item Multitasking (up to 56 tasks)
  \item Preemptive:  64 unique priority levels.
  Always runs highest priority task.
  \item Deterministic:  ISRs etc. take known,
  fixed amount of time.
  \item Individual stacks for each task. Can
  be different sizes.
  \item Services:  mailboxes, queues, semphores,
  dynamic memory allocation.
  \item Interrupt management.  Scheduler can
  check for a task that should be activated after an
  interrupt.
  \item Mature.  In use since 1992.  Incorporated
  into commercial products.
\end{itemize}

%\end{slide}
%\begin{slide}


%%%%** Section 2 
\section{$\mu$C/OS-II   Tasks and Calls}

%%%%** Section 2.1
\subsection{Task Priorities}

All tasks have a unique priority.	%<*chn>
Low priority numbers
have the  highest priority.
Priority numbers 0,1,2,3, are reserved.

There are a
{\it maximum } of 64 total priority levels, but if you need
smaller memory useage, the system can be configured for
fewer.  Thus the lowest priority task is not always 64. Instead
the OS defines four levels at the bottom:

{\tt OS\_LOWEST\_PRIO, OS\_LOWEST\_PRIO-1, OS\_LOWEST\_PRIO-2, OS\_LOWEST\_PRIO-3 }

(which would be 64, 63 ,62, and 61 in a full system).  The lowest user level priority task is therefore:
{\tt OS\_LOWEST\_PRIO - 4 }.   The four
highest priorities are also
all reserved.
The maximum number of priorities
is thus 56.

Only one task may have each priority level so 56 is also the maximum
number of tasks.  Since task priorities are unique they are also
used to identify the processes (like the pid in Unix).
The max number of tasks can be reduced at compile time
to save memory.

%
%\begin{itemize}
%        \item All tasks have a unique priority.
%        \item Priority values can be used to ID tasks.
%        \item 4 Highest Priority numbers (0,1,2,3) are reserved.
%
%        \item 4 Lowest Priority numbers are reserved:
%               {\tt OS\_LOWEST\_PRIO, OS\_LOWEST\_PRIO-1, OS\_LOWEST\_PRIO-2, OS\_LOWEST\_PRIO-3 }
%
%        \item Lowest Usable Priority available to users:
%               {\tt OS\_LOWEST\_PRIO - 4 }
%        \item Normally {\tt OS\_LOWEST\_PRIO = 64}  but can configure
%        system to fewer tasks for smaller memory use.
%\end{itemize}
%
	%<*>
%%%%** Section 2.2
\subsection{OS Function Calls (selected)}

 Tables \ref{Ta:FC0} and \ref{Ta:FC1}	%<chn>
 summarize key function calls in $\mu$C/OS-II.	%<chn>

% See Tables from Notes!!   too big for slides.

	%<*chn>
%%%%** Table 1 
\begin{table}
\begin{center}
 %p = paragraph, c = center, l = left, r = right
  \begin{tabular}{|l|p{1.6in}|p{2.75in} |}
  \hline
Name   &  Prototype  &  Function \\ \hline
 OSInit()       & {\tt void OSInit(void) }      & Sets up OS internal data
                                                  structures. Must be called first. \\ \hline
 OSStart()      & {\tt void OSStart(void) }     & Begins operation of the scheduler.\\ \hline
 OSIntEnter()   & {\tt void OSIntEnter(void)}   & Must be called at start of an ISR.
                                                  Lets OS keep track of number of ISRs. \\ \hline
 OSIntExit()    & {\tt void OSIntExit(void)}    & Call at end of an ISR
%                                              {\it before} restoring conext and {\tt IRET}.
                                                  \\ \hline
 OSTaskCreate() & {\tt INT8U OSTaskCreate( void (*task)(void *pd), void *pdata, OS\_STK *ptos, INT8U prio) }
                & Creates a new task.
                 {\tt task} is a pointer to the task's function,
                 {\tt pdata} points to optional parameters to the task,
                 {\tt ptos} is a pointer to the tasks stack top.
                 {\tt prio} is the task priority.  Note that since each
                 $\mu$COS-II task must have a unique priority, priority is
                 sometimes used to as the ``task number'' as well (see
                 OSTaskDel().
                   \\ \hline
OSTaskDel()     & {\tt INT8U OSTaskDel(INT8U prio) }    &
                Deletes a task. The task is indicated by its priority
                number so any task can delete any other task.
                {\tt OS\_PRIO\_SELF} is a predefined constant
                which a task can use to delete itself if it doesn't know
                its priority. \\ \hline
OSTaskDelReq() & {\tt INT8U OSTaskDelReq(INT8U prio) }  &
                Send a request to a task to delete itself.   If a task is
                set up for this, it will {\it also} issue a
                OSTaskDelReq(OS\_PRIO\_SELF) and if that returns
                OS\_TASK\_DEL\_REQ, then it should delete itself.  \\ \hline
OSTaskSuspend() & {\tt OSTaskSuspend(INT8U prio)}       &
                Block the execution of the current task.    \\ \hline
OSTaskResume()  & {\tt OSTaskResume(INT8U prio)}        &
                Resume execution of a task which was blocked using
                OSTaskSuspend()  \\ \hline
OSTaskChangePrio()& {\tt INT8U OSTaskChangePrio( INT8U oldprio, INT8U newprio)} &
                Change the priority of a task.  {\tt oldprio } must exist
                and {\tt newprio } must be free.    \\ \hline
OSSchedLock()   & {\tt void OSSchedLock(void)} &
                Stops all task scheduling but not interrupts.
                Higher priority tasks will no longer take over. \\ \hline
OSSchedUnlock() & {\tt void OSSchedUnlock(void)} &
                Resumes task scheduling. \\ \hline
OSTimeDly()     & {\tt void OSTimeDly(INT16U ticks)} &
                Delay the calling task by {\tt ticks} clock ticks (i.e. timer0
                interrupts).  The scheduler will start other tasks during
                this delay       \\ \hline
OSTimeDlyHMSM() & {\tt void OSTimeDlyHMSM( INT8U hours, INT8U minutes, INT8U seconds, INT8U milli) }    &
                Same as ISTimeDly() but uses natural time units.         \\ \hline
OSTimeDlyResume()       & {\tt INT8U OSTimeDlyResume(INT8U prio) }      &
                Wakes up a task which has called either of the two functions
                above.   \\ \hline
OSTimeGet()     & {\tt INT32U OSTimeGet(void)}  &
                Returns the current time measured by number of ticks.
                The result is a 32 bit unsigned integer.  \\ \hline

\end{tabular}
\caption{Table of Selected $\mu$C/OS-II calls}
\label{Ta:FC0}
\end{center}
\end{table}

%%%%** Table 2 
\begin{table}
\begin{center}
 %p = paragraph, c = center, l = left, r = right
  \begin{tabular}{|l|p{1.6in}|p{2.75in} |}
  \hline
Name   &  Prototype  &  Function \\ \hline
OS\_ENTER\_CRITICAL()   & {\tt }        &
                Not actually a function.  This {\it macro} inserts
                the machine instruction into
                your code to block all interrupts .      \\ \hline
OS\_EXIT\_CRITICAL()    & {\tt }        &
                Not actually a function.  This {\it macro} inserts
                the machine instruction to enable interrupts. \\ \hline
OSSemCreate()   & {\tt OS\_EVENT *OSSemCreate( WORD value)}     &
                Create a new semaphore and initialize it to
                {\tt value}. Returns a pointer to the semphore
                (actually an {\it event control block}).  \\ \hline
OSSemPost()     & {\tt INT8U OSSemPost(OS\_EVENT *pevent }      &
                Signal the semaphore pointed to by {\tt pevent}.
                (enables other tasks to continue).   \\ \hline
OSSemPend()     & {\tt void OSSemPend(OS\_EVENT *pevent, INT16U timeout, INT8U *err }   &
                Wait for access to resource controlled by semaphore {\tt pevent}.
                {\tt pevent} must have been previously created by
                OSSemCreate(). {\tt timeout} is a time value (in ticks)
                which allows the task to wake up  if the semaphore
                never becomes available.  If this happens, the
                function will return {\tt OS\_TIMEOUT}. Do not use inside
                an ISR.   \\ \hline

 \end{tabular}

\caption{Table of Selected $\mu$C/OS-II calls (cont.)}
\label{Ta:FC1}
\end{center}
\end{table}


	%<*>
%\end{slide}

%\begin{slide}

\newpage	%<chn>
%%%%** Section 3 
\section{Code Example}
%\tiny{
\begin{verbatim}
//----------------------------------------------------------------------------
// AUTHOR:    Jered Aasheim
// CREATED:   February 23, 2001
// MODIFIED:  February 23, 2001
// PROJECT:   EE472 Final Project
// MODULE:    Test Application using uC/OS RTOS
//
// Abstract:  The following program demonstrates how to use the uC/OS Real-Time
//            Operating System to solve a critical section problem.  Using the
//            uC/OS facilities, two different tasks are created and a binary
//            semaphore is used to protect a critical section (in this case the
//            first row on the LCD screen).
//
//            Unfortunately, this program only demonstrates a few facilities of
//            the uC/OS RTOS.  For further examples using uC/OS, refer to
//            "An Embedded Software Primer" by David E. Simon.
//
//            NOTE: In order to use the uC/OS RTOS, you will need to add the
//                  UCOS.LIB and UCOS.H files to your Paradigm C++ project.
//
//----------------------------------------------------------------------------

#include "td40.h"
#include "ucos.h"

#define STACK_SIZE              1024
#define WAIT_FOREVER            0


#define STARTUP_PRIORITY         0
#define TASK1_PRIORITY          11
#define TASK2_PRIORITY          12

UWORD StartUp_Stk[STACK_SIZE];
UWORD T1_Stk[STACK_SIZE];
UWORD T2_Stk[STACK_SIZE];

OS_EVENT *pBinSem = NULL;               // Binary semaphore

static UBYTE err = 0;

void SetupTimers(void);

void far StartUpTask(void* data);

void far SimpleTask0(void* data);
void far SimpleTask1(void* data);

void interrupt far Timer0_ISR(void);



void main(void)
{
   ae_init();
   lcd_init();

   OSInit();                            // Initialize uC/OS

   //SetupTimers();                     // Configure the TD40 timers.
   setvect(uCOS, (void interrupt (*)(void))OSCtxSw);

   pBinSem = OSSemCreate(1);            // Initialize the semaphore

   OSTaskCreate(StartUpTask, (void*)0, (void*)&StartUp_Stk[STACK_SIZE-1], STARTUP_PRIORITY);

   OSStart();                           // Start the uC/OS
}


void far StartUpTask(void* data)
{
        OS_ENTER_CRITICAL();

   // Install the uC/OS timer ISR
   setvect(0x0008, (void interrupt (*)(void))OSTickISR);
   // Install applications timer ISR
   setvect(0x0081, Timer0_ISR);

   // Initialize Timer0
   outport(0xFF32, 0x0007);         // Unmask TIMER0 interrupt and set Timer0
                                                                                        // interrupt to low priority

   outport(0xFF56, 4000);                       // Disable TIMER0

   outport(0xFF52, 10000);          // Initialize TIMER0's count

   outport(0xFF56, 0xE001);         // Enable TIMER0 to generate an interrupt,
                                    // use maxcount A, don't use prescaled value,
                                    // use internal clock, continuous operation

        OS_EXIT_CRITICAL();

   OSTaskCreate(SimpleTask0, (void*)'1', (void*)&T1_Stk[STACK_SIZE-1], TASK1_PRIORITY);
   OSTaskCreate(SimpleTask1, (void*)'2', (void*)&T2_Stk[STACK_SIZE-1], TASK2_PRIORITY);
   OSTaskDel(OS_PRIO_SELF);

}


//----------------------------------------------------------------------------
// SimpleTask() - demonstrates a simple task that uses semaphores.  In
//                this case, the LCD screen is the critical section.
//----------------------------------------------------------------------------
void far SimpleTask0(void* data)
{
   for(;;)
   {
      // Get the semaphore (i.e. P())
      OSSemPend(pBinSem, WAIT_FOREVER, &err);

      // Critical section...
      lcd_movecursor(0x80);
      lcd_putstr("Task ");
      lcd_put((char)data);
      delay_ms(2000);

      // Give up the semaphore (i.e. V())
      OSSemPost(pBinSem);

      // Block for 5 ticks - forces a task switch
      OSTimeDly(5);
   }
}


void far SimpleTask1(void* data)
{
   for(;;)
   {
      // Get the semaphore (i.e. P())
      OSSemPend(pBinSem, WAIT_FOREVER, &err);

      // Critical section...
      lcd_movecursor(0x80);
      lcd_putstr("Task ");
      lcd_put((char)data);
      delay_ms(2000);

      // Give up the semaphore (i.e. V())
      OSSemPost(pBinSem);

      // Block for 5 ticks - forces a task switch
      OSTimeDly(5);
   }
}



//----------------------------------------------------------------------------
// SetupTimers() - initializes the timers on the TD40 board.
//                 The uC/OS requires a timer to periodically
//                 generate an interrupt so that it can schedule
//                 tasks appropriately.  In most systems, however,
//                 uC/OS can't just define the ISR for a Timer
//                 for several reasons:
//
//                 1) Defining a timer ISR is platform specific
//                    and uC/OS is intended to be platform independent.
//
//                 2) The application may need to define its own
//                                                       ISR for the timer.  If uC/OS defines the ISR
//                    then the application can't define their own
//                    ISR.
//
//                 The solution is simple: have the application
//                 configure the timer to generate interrupts and
//                 the uC/OS "timer" ISR.  Then, inside the uC/OS
//                 ISR a software interrupt is generated (i.e. INT 81)
//                 which runs the users "timer" ISR.  Using this
//                 technique, after all is said and done, uC/OS will
//                 have the required timer interrupt ISR setup correctly
//                 and the application can define its own timer ISR.
//
//                 NOTE: Even if the application doesn't use timer interrupts
//                       it still must define a "dummy" timer ISR (see below).
//
//----------------------------------------------------------------------------
void SetupTimers(void)
{
   // Install the uC/OS "task switch" ISR
   //setvect(uCOS, (void interrupt (*)(void))OSCtxSw);

   // Install the uC/OS timer ISR
   setvect(0x0008, (void interrupt (*)(void))OSTickISR);

   // Install applications timer ISR
   setvect(0x0081, Timer0_ISR);

   // Initialize Timer0
   outport(0xFF32, 0x0007);               // Unmask TIMER0 interrupt and set Timer0
                                          // interrupt to low priority

   outport(0xFF56, 4000);                 // Disable TIMER0

   outport(0xFF52, 10000);                // Initialize TIMER0's count

   outport(0xFF56, 0xE001);               // Enable TIMER0 to generate an interrupt,
                                          // use maxcount A, don't use prescaled value,
                                          // use internal clock, continuous operation
}



//----------------------------------------------------------------------------
// Timer0_ISR() - "dummy" Timer0 ISR, see SetupTimers() above
//                for details
//----------------------------------------------------------------------------
void far interrupt Timer0_ISR()
{
   outport(0xFF22,0x0008);          // Non-specific EOI
}
\end{verbatim}
%}    % end of tiny

%\end{slide}


%\end{document}

